% !TEX root = ../journal_v2_0421052099.tex
%
% Condensed from the discussion and threats to validity of
% thesis_book/chapters/evaluation.tex.
\section{Discussion}\label{sec:discussion}

\subsection{Findings against the Research Questions}

\textbf{RQ1.} Agentic navigation improves multi-hop accuracy, and the multi-hop
qualifier is doing real work. AGR leads every accuracy column, and all eight
paired tests reject. On CWQ its accuracy rises with hop count, while every
static system decays. Three boundaries come with that answer. Think-on-Graph
hedges marginally less on CWQ, at $22.5\%$ against $23.0\%$, the one cell of
\Cref{tab:main} that AGR does not lead. It is also better on CWQ's one-hop
stratum. And the margin over it is located in the budget, against a baseline
that prunes from narrower candidate sets. So the claim defended here is that
guided search completes under a realistic budget where beam search does not.

\textbf{RQ2.} The answer has three parts that a yes-or-no question would merge.
First, navigation eliminates structural hallucination, and the verifier is not
what does it. Second, the precision advantage over Think-on-Graph belongs to the
architecture, and this experiment cannot award any of it to the layer. Removing
the verifier moves per-question precision by $+0.001$ and $-0.010$, an order of
magnitude below the gap. Entity filtering runs only on the give-up path, which
$10$ WebQSP and $29$ CWQ questions entered. A filter that touches $2.5\%$ of a
test set cannot produce a $0.13$ precision margin on it. Two differences that
the ablation holds fixed remain live. One is the groundedness filter on
evaluator resolutions, which is navigation rather than verification. The other
is the drafting prompt's provisions, which the baselines' plain prompt lacks.
Third, verification shows no measurable Hits@1 effect. A metric satisfied by one
correct entity among several cannot see a mechanism that works on the others.
What the layer offers is an answer that arrives with the triples supporting it.
That is an auditability claim, not an accuracy claim.

\textbf{RQ3.} One condition of four reaches significance. Removing the planner
improves WebQSP by $0.083$ F1 at $p = 0.006$ while saving $31\%$ of tokens, and
it costs CWQ $0.047$ F1 at $p = 0.088$. That asymmetry, not a pooled average, is
the result. It says how decomposition should be applied, not whether it works.
Backtracking, verification, and the scoring blend show no detectable effect at
$n \approx 200$. On cost, the whole apparatus runs at half the calls of the
agentic baseline.

So the three answers do not point the same way. Navigation is where the accuracy
comes from, and verification is where the auditability comes from.
Decomposition, the component the literature treats as straightforwardly
beneficial, is the one this work can show is sometimes harmful.

\subsection{What Agency Buys, and What It Costs}

Agency buys the deep tail. On CWQ the agentic system's accuracy rises across hop
strata while every static system decays. It buys that at roughly an order of
magnitude more tokens than one-shot retrieval. Against the other agentic system,
it buys higher accuracy for half the calls, and it never exhausts the call
budget that clips that baseline on $44\%$ of CWQ questions. That is the
difference between a guided search and a blind one, priced. Agency also
lengthens the failure surface. Decomposition errors, abandoned relations, and
verifier errors are failures that a static baseline cannot have. It has no plan
to mis-decompose, no frontier to abandon, and no draft to check. Agency converts
a retrieval problem into a control problem, and control problems fail in more
ways.

\subsection{When Verification Helps and When It
\texorpdfstring{\mbox{Over-Hedges}}{Over-Hedges}}

The layer's contribution is auditability, and its cost is recall. Both
directions show up at case level, which is where they are established, since the
aggregate effect of removing the layer is indistinguishable from zero. The layer
helps when a draft asserts something the traversal never supported. It rejected
the fabricated Brando claim of \Cref{sec:errors} before it could be emitted. A
pair of sibling questions about the battle of Vicksburg shows the other kind of
help. One asks who was president during the battle, and its gold answer is
Ulysses S.~Grant, who took office six years later. AGR answered that the claim
could not be verified, and its verifier flagged the draft. The ablated system
without a planner asserted Grant with confidence and was scored a hit for
restating a gold value that answers a different question. The system that hedged
was the more correct one, and the metric scored it as a loss.

The layer over-hedges when a true claim has no single edge stating it in the
drafted phrasing. The five wrongly rejected cases of \Cref{sec:errors} show
that, and four of them were retries that ran and lost the answer. So the claim
is narrow. Verification does not make this system more accurate. It makes the
answers it does give attributable, at a small and measurable cost in recall. And
it provides no protection against an answer that is grounded, connected, true,
and not what was asked.

\subsection{Threats to Validity}\label{sec:threats}

\textbf{The environment is friendlier than the knowledge base it stands in for.}
It is a union of subgraphs that were extracted to contain their own answers.
That is why reachability is certified at $97.3\%$ and $99.7\%$, and why the
accuracies are not comparable with evaluations over full Freebase. Every system
navigates the identical environment, so the relative ordering is sound even
where the absolute level is optimistic. Every claim in this paper is relative
for that reason.

\textbf{An identical environment is not identical access.} The agentic baseline
sees narrower candidate sets than AGR. The static baseline sees a single logical
hop under an unordered fanout cap. So the reported margins are bounded in both
directions, and equal-width re-runs would settle them.

\textbf{Every number comes from one run per condition.} The intervals and paired
tests resample questions, so they bound the sampling variance of the draw and
nothing else. The backbone repeated its trajectory on only $8$ of $12$ probes.
The margins defended here are large relative to their intervals, so a rank
reversal from decoding noise alone is unlikely. But the experiment that would
bound that variance was not run.

\textbf{Entities are keyed by name.} Homonyms collapse into one node. A merged
node can manufacture a path between unconnected things, which inflates
reachability and lets a wrong answer verify as grounded. So structural
groundedness is a claim about this graph's node identity, not about the world.

\textbf{Wrongful acceptance is unmeasured.} The verifier logs the claims it
rejects and not the ones it accepts. So one polarity of its error is reported at
census scale and the other from a single specimen. The exposed population is
bounded only as $[39, 2{,}008]$ of $2{,}008$ accepted claims.

\textbf{Further limits.} The judge fell short of its pre-specified bar by
$0.0005$ and shares the systems' backbone. Topic entities are given rather than
linked, so the accuracies assume a linking step that a deployed system would
perform itself. The ablations are underpowered. The extraction bug makes the
measured accuracy a lower bound on one question shape. And the benchmark defects
were adjudicated by the first author, with the consensus signal visible.

\textbf{One criterion is not reported.} Path fidelity asks whether a system
reaches the answer by the relation chain that the benchmark's own query
specifies. It needs the gold SPARQL chains, and the distribution used here
publishes triples without them. Reconstructing the chains would measure whether
the right answer came by the right path. The groundedness metrics can only
approximate that question, and it bears directly on the echo attractor.

\endinput
